\documentclass[10pt,a4paper]{amsart}
\usepackage[margin=19mm]{geometry}
\usepackage{amsmath,amssymb,mathtools}
\usepackage[scr=rsfs,cal=euler]{mathalfa}
\usepackage{xfrac}
\usepackage[unicode,hidelinks,bookmarks=false]{hyperref}
\newcommand{\ie}{i.e.\ }
\newcommand{\cf}{cf.\ }
\newcommand{\resp}{respectively\ }
\newcommand{\Subsection}{\subsection}
\providecommand{\nobreakdash}{\nobreak\hbox{-}}
\setlength{\emergencystretch}{2em}
\multlinegap0pt
\usepackage[shortlabels]{enumitem}
\usepackage{amssymb}
\usepackage{mathtools}
\newtheorem{lemma}{Lemma}
\newcommand{\calH}{\mathcal H}
\title{A sparse transference principle \\for a non-monotone Ramsey property}
\author{Gaia Carenini}
\date{Source: arXiv:2606.02873v1 (CC BY 4.0)}
\begin{document}
\maketitle

\noindent\emph{Source and licence.} A sparse transference principle \\for a non-monotone Ramsey property. Version arXiv:2606.02873v1 (CC BY 4.0), distributed by arXiv under the Creative Commons Attribution 4.0 International licence, http://creativecommons.org/licenses/by/4.0/, as recorded in the arXiv licence metadata for this exact version and shown on the abstract page https://arxiv.org/abs/2606.02873v1. Full licence notice: \texttt{licenses/arxiv-2606.02873-CC-BY-4.0.txt}.\par\smallskip

\noindent\textit{Editorial scope.} Counted unit: the article's lemma lem:codegrees (source0.tex lines 431--441). The article's complete original proof of it is delivered with it (source0.tex lines 443--465, one \texttt{proof} environment of 23 lines). No mathematics is written, repaired or re-derived: the statement and the proof are the article's own lines, in the article's own notation, and no retained line is substituted or re-flowed.  No further source-local context is needed: every cross-reference of every retained line resolves inside the two delivered spans.  Nothing was omitted from any retained span, so the package carries no omission record.  No retained line contains a citation, so the wrapper carries no bibliography and no bibliography entry.  No retained line contains a figure environment, an image-inclusion command or drawing code, so no image is delivered and the package declares no asset dependency.  Editorial formatting only, in addition: an AMS \texttt{amsart} preamble that restates the article's own declarations and macros used by the retained lines, namely the environments \texttt{lemma} on separate counters, together with the standard packages those lines need.  The retained environments print in this document as Lemma 1 in order of appearance, whereas the article prints the same items with the numbering of its own full document.  The complete native source of the article as distributed in the arXiv e-print for this exact version is bundled unchanged as \texttt{sources/201934/source0.tex}.
\begin{lemma}[Codegrees in a dense random host]
\label{lem:codegrees}
With high probability, $F\sim G(N,1/2)$ satisfies:
\begin{enumerate}[label=\textup{(\arabic*)}]
\item $|E(F)|=\Theta(N^2)$;
\item $e(\calH_F)=\Theta_H(N^h)$;
\item the average degree of $\calH_F$ is $d=\Theta_H(N^{h-2})$;
\item if $\tau=N^{-1/m_2(H)}$, then, there exists a constant $A_H$,
$ \Delta_\ell(\calH_F)\le A_H d\tau^{\ell-1}$ for every $2\le \ell\le q$.
\end{enumerate}
\end{lemma}

\begin{proof}
The first two assertions follow from the usual concentration estimates for the
number of edges and for the number of induced copies of a fixed graph in
$G(N,1/2)$.  They imply the third, since
$$
        d=\frac{qe(\calH_F)}{|E(F)|}=\Theta_H(N^{h-2}).
$$

For the codegree estimate, fix a set $\sigma\subseteq E(F)$ of $\ell$ host
edges.  If $\sigma$ is contained in no induced copy of $H$, then its
codegree is zero.  Otherwise, in every induced copy containing $\sigma$, the
edges of $\sigma$ correspond to an $\ell$-edge subgraph $J\subseteq H$.
For a fixed such $J$, once the images of the vertices incident with these
$\ell$ edges are fixed, the remaining vertices of the copy can be chosen in
at most $O_H(N^{h-v(J)})$ ways.  Hence
$$
        \Delta_\ell(\calH_F)
        \le C_H\max_{\substack{J\subseteq H\\ e(J)=\ell}}N^{h-v(J)}.
$$
Since $\ell\ge2$, every contributing $J$ has $v(J)\ge3$.  By the
definition of $m_2(H)$, $\ell-1\le m_2(H)(v(J)-2)$,
and therefore $ N^{h-v(J)}\le N^{h-2-(\ell-1)/m_2(H)}$. Together with $d=\Theta_H(N^{h-2})$, this gives the required bound.
\end{proof}

\end{document}
